\scanpage{98}
On sait que cette dernière est \unsure{équivalente} à la cat. ponctuelle. Donc $\mathrm{DR}_{\mathrm{sp}}(X,k)_0$ est équivalente à la cat. ponctuelle.

Si maintenant on travaille dans $\mathrm{DR}_{\unsure{\mathrm{spufl}|\mathrm{spu\,fl}}}$ \deleted{(\unsure{resp.|cat.} $\mathrm{DR}_{\unsure{\mathrm{spufl}|\mathrm{spfl}}}$)} i.e. avec des $\unsure{A|\mathcal A}^{\unsure{**|\alpha\alpha}}$ \inserted{u.} flasques, ($A^{ij}$ \inserted{u.} flasque si $0<j$) alors $\Phi$ est \inserted{u.} flasque. \deleted{(\unsure{resp. univ. flasque}).}
Supposons que ceci \unsure{implique} que les $\Gamma_k^{i}\Phi$ $(i>0)$ sont \unsure{univ.} flasques, alors $\mathrm{DR}(\Phi,T)$ est aussi flasque, \deleted{\unsure{et}} et $\unsure{=}$ univ. flasque, et on voit comme dessus que la cat. dérivée est \unsure{équiv.|équivalente} à celle construite sur \deleted{la cat. des} $\mathrm{DR}(\Phi,T)$, avec $\Phi$ univ. flasque et $\Phi/k.T$ plat, i.e. avec les \unsure{suite|suites} exactes $0\unsure{\longrightarrow|\text{---}} k\longrightarrow\Phi\longrightarrow\Psi\longrightarrow0$, avec $\Phi$ univ. flasque et $\Psi$ $k$-plat ; or cette cat. dérivée est encore la catégorie ponctuelle. On trouve donc

\uline{Prop} La cat. $\mathrm{DR}_{\mathrm{sp}}(X,k)_0$ est équivalente à la cat. ponctuelle. Il en est de $=$ de la cat. $\mathrm{DR}_{\mathrm{sp\,ufl}}(X,k)_0$, si on suppose que pour \unsure{tt} \inserted{$k$-Module} $\Phi$ plat et univ. flasque, les $\Gamma^i\Phi$ $(i>0)$ sont aussi univ. flasques.
% archive sheet number: 97
